% =============================================================================
%  06-placement.tex -- Placement integration: how the node tree binds to the
%  placement layer (value types, per-edge math, resolver, diagnostics gate).
% =============================================================================

\section{Placement integration}\label{sec:placement}

The node tree of \cref{sec:tree} stores structure and intent: which part attaches under which, and
one \code{StationLink} per edge describing \emph{how} the child seats against its parent. Nothing in
the stored model is an absolute coordinate. Everything spatial --- each part's pose in the rocket
frame, the composite center of mass, the overlap verdict, the visualizer's mesh positions --- is
derived on demand from that intent by the placement layer, a small vocabulary of value types and
free functions in \srcf{core/model/parts/Placement.h}. The layer is pure geometry: no time, no mass,
no tree access. The tree walk that drives it lives with the tree (\cref{sec:placement-resolver}),
and its output is cached on the node behind the placement gate of \cref{sec:caching}.

This section reads bottom-up: the value types first, then the per-edge placement math, then the
node-tree resolver that composes edges into absolute poses, and finally the diagnostics --- the
per-seam radial check and the whole-tree envelope sweep --- whose verdict gates every composite
read.

\subsection{The placement vocabulary}\label{sec:placement-vocab}

\Cref{fig:placement-types} shows the complete layer. Every type is a plain value; every function is
free. The layer deliberately owns nothing and observes parts only through \code{const} references
and non-owning pointers, so it composes with the single-ownership story of \cref{sec:tree} without
touching it.

\classfig{class-placement-types}{The placement layer of
\srcfcap{core/model/parts/Placement.h}: stored intent (\code{StationLink}), derived geometry
(\code{Station}, \code{Pose}, \code{Placed}), diagnostics values (\code{OverlapDiagnostic},
\code{SolveResult}), and the free functions that connect them.}{fig:placement-types}{0.95}

\paragraph{StationLink: the stored intent.}
A \code{StationLink} (\src{core/model/parts/Placement.h}{54}, \cref{lst:placement-stationlink}) is
the only spatial relationship the model persists: a fractional station on the parent, a fractional
station on the child, a gap in metres, a seat kind, and a per-seam orientation quaternion. Stations
are fractions of \emph{live} length --- $0$ names the aft plane, $1$ the fore plane --- and are
resolved against \code{getLength()} on every read by \code{Part::stationAt}
(\src{core/model/parts/Part.cpp}{44}), which clamps the fraction and maps it to
$z = (\mathit{station01} - 1) \cdot L$ in the part's local \zfwd{} frame.

\begin{listing}[htbp]
\begin{minted}{cpp}
struct StationLink
{
    double     parentStation01{0.0};             // 0 = aft plane, 1 = fore plane, on the PARENT
    double     childStation01{1.0};              // 0 = aft plane, 1 = fore plane, on the CHILD
    double     gap{0.0};                         // meters; meaning per SeatKind
    SeatKind   seat{SeatKind::Abut};             // selects the radial check and the gap sign
    Quaternion childRot{Quaternion::Identity()}; // (x,y,z,w) seam orientation; identity in 3-DOF
};
\end{minted}
\caption{\code{StationLink} (\srccap{core/model/parts/Placement.h}{54}). The default-constructed
link means ``child fore plane abuts parent aft plane'' --- the natural nose $\to$ body $\to$ tail
build order, so a design assembled with default links stacks aft without any authored
numbers.}\label{lst:placement-stationlink}
\end{listing}

\begin{keyinsight}[Length edits move the seam]
Because stations are stored as fractions and resolved against \code{getLength()} at read time, a
link never goes stale. Stretch a body tube and every seam authored against its planes --- and every
part stacked behind it --- moves with it on the next resolve, with no fix-up pass and no stored
coordinate to forget to update.
\end{keyinsight}

\begin{rejected}[Rejected --- authored absolute offsets]
The alternative is to store each child's axial offset directly, either from the root or from its
parent's origin. That representation is simpler to resolve but wrong under edit: an authored
coordinate is a snapshot of geometry that was true when it was written, and every length change
afterwards silently invalidates it. Fractional stations keep the stored quantity in intent space
(``my fore plane against your aft plane''), where it remains true under any resize.
\end{rejected}

\paragraph{Seat kinds.}
\code{SeatKind} (\src{core/model/parts/Placement.h}{32}) is a closed three-value enum, because a
rigid axisymmetric stack admits exactly three radial relationships between mated parts. The seat
does double duty: it selects which radial compatibility check applies at the seam
(\cref{sec:placement-diagnostics}), and it fixes the \emph{sign} of the gap so the author always
writes a non-negative number (\cref{tab:placement-seats}). An arbitrary axial position needs no
fourth kind --- it is already expressible as a station fraction.

\begin{table}[htbp]
\centering\small
\begin{tabularx}{\linewidth}{@{}llXX@{}}
\toprule
Seat & Seam geometry & Radial check & Gap meaning (sign in \zfwd) \\
\midrule
\seat{Abut} & rim to rim & outer radii equal & forward standoff ($+z$) \\
\seat{NestInBore} & child inside parent bore & child OD $\le$ parent bore & insertion depth ($-z$, child moves aft) \\
\seat{OnSurface} & child on parent outer wall & outer radii equal & axial standoff ($+z$) \\
\bottomrule
\end{tabularx}
\caption{The three seat kinds. Each fixes both the seam's radial compatibility test and the
direction a positive gap displaces the child, so no authored negative gaps
exist.}\label{tab:placement-seats}
\end{table}

Three snap verbs --- \code{abut(gap)}, \code{nestInBore(depth)}, and
\code{seatOnWall(parentStation01)} (\src{core/model/parts/Placement.h}{134}) --- produce a
\code{StationLink} by value for the common cases. They are pure sugar over the brace-initializer
form: no pose is stored and no part is mutated, they simply read better as verbs at call sites such
as the CLI's link parser (\src{cli/Repl.cpp}{226}), which seeds each seat's canonical station pair
from its verb before applying the author's explicit overrides (\cref{sec:cli}).

\paragraph{Reserved orientation.}
\code{childRot} is carried in every link but is the identity throughout the current 3-DOF regime.
It exists so the placement schema needs no change when 6-DOF arrives: \code{Pose::compose} already
multiplies orientations, and a stored 3-DOF link is bit-identical to a pure-translation tree. The
cost of this reserve is one identity quaternion multiply per child per resolve.

\paragraph{Derived geometry.}
A \code{Station} (\src{core/model/parts/Placement.h}{67}) is a resolved axial landmark --- $z$,
outer radius, inner radius --- produced by \code{stationAt} and never stored. A \code{Pose}
(\src{core/model/parts/Placement.h}{79}) is a rigid transform into the sub-tree root frame: a
fore-plane origin on the axis plus an orientation, with \code{compose} implementing the
rotate-then-translate composition law (vector addition while orientations are identity). A
\code{Placed} (\src{core/model/parts/Placement.h}{98}) pairs a non-owning part pointer with its
resolved pose and its seat parent's id. The resolver emits \code{Placed} entries in deterministic
depth-first attachment order, parent before child, which yields two guarantees consumers lean on:
the envelope sweep is reproducible, and \code{parentId} always names an \emph{earlier} entry in the
vector (0 for the sub-tree root), so a consumer can reconstruct the tree relation by streaming the
vector once, without access to the ownership tree itself.

\paragraph{Diagnostics values.}
An \code{OverlapDiagnostic} (\src{core/model/parts/Placement.h}{109}) locates one violation:
offender id, host id, the worst station, the penetration depth, and a human-readable message. Parts
are named by stable id (\code{Part::getId()}), never by name --- names are not unique --- and never
by pointer, which would dangle if the tree were edited between resolve and report. A
\code{SolveResult} (\src{core/model/parts/Placement.h}{123}) bundles the verdict flag with the
diagnostics list; \cref{sec:placement-gate} covers how that verdict is enforced.

\subsection{Per-edge geometry: placeChild}\label{sec:placement-placechild}

All per-edge placement math is one pure function, \code{placeChild}
(\src{core/model/parts/Placement.cpp}{39}, \cref{lst:placement-placechild}). Given the parent's
already-resolved pose, the two parts, and the link binding them, it asks each part for the live
landmark at its chosen fractional station, lines the two stations up, and displaces the child by
the seat-signed gap.

\begin{listing}[htbp]
\begin{minted}{cpp}
Pose placeChild(const Pose& parentPose, const Part& parent, const Part& child,
                     const StationLink& link)
{
    const Station p = parent.stationAt(link.parentStation01);  // landmark on the parent
    const Station c = child.stationAt(link.childStation01);    // landmark on the child

    // signedGap encodes seat direction in +z = forward: NestInBore inserts aft (-z, gap is depth),
    // Abut/OnSurface stand off forward (+z). The author always writes a non-negative number.
    const double signedGap = (link.seat == SeatKind::NestInBore) ? -link.gap : link.gap;

    // Line the child's station up with the parent's, then displace by the gap. The pose locates the
    // child's fore plane, which sits c.z from its chosen station -- hence the `- c.z`.
    const double childOriginZ = p.z + signedGap - c.z;

    const Pose childInParent{Vector3(0.0, 0.0, childOriginZ), link.childRot};  // coaxial in 3-DOF
    return parentPose.compose(childInParent);
}
\end{minted}
\caption{The complete per-edge placement math
(\srccap{core/model/parts/Placement.cpp}{39}).}\label{lst:placement-placechild}
\end{listing}

The central line reads as physical intent: \cppinline|childOriginZ = p.z + signedGap - c.z|. The
parent's chosen station sits at \code{p.z} in the parent frame; the child must be positioned so its
own chosen station lands there; and because a pose locates the child's fore plane --- its local
origin --- rather than the chosen station, the station's own offset \code{c.z} is subtracted back
out. The seat supplies the gap's direction, so \seat{Abut} and \seat{OnSurface} standoffs move the
child forward while a \seat{NestInBore} depth drives it aft into the bore. In 3-DOF the child is
coaxial with its parent ($x = y = 0$); radial position is a property the diagnostics \emph{check},
not one the author places.

\subsection{The node-tree resolver}\label{sec:placement-resolver}

Turning a tree of links into absolute poses is a depth-first walk, and the walk lives with the tree:
\code{resolveNodes} (\src{core/model/PartsModel.cpp}{29}, \cref{lst:placement-resolvenodes}) is a
file-local function of \srcf{core/model/PartsModel.cpp} that plants the sub-tree root at the local
origin (\code{Pose\{\}}) and composes one \code{placeChild} result per edge on the way down,
emitting one \code{Placed} per node in attachment order.

\begin{listing}[htbp]
\begin{minted}{cpp}
/// Depth-first resolve of a node tree from @p pose: one Placed per node in attachment order,
/// parent before child.
void resolveNodes(const PartNode& n, const part::Pose& pose, part::PartId parentId,
                  std::vector<part::Placed>& out)
{
    out.push_back(part::Placed{&n.part(), pose, parentId});
    for(const auto& child : n.children())
    {
        resolveNodes(*child, part::placeChild(pose, n.part(), child->part(), child->link()),
                     n.part().getId(), out);
    }
}
\end{minted}
\caption{The node-tree resolver (\srccap{core/model/PartsModel.cpp}{29}). The tree walk is ten
lines because all geometry lives in \code{placeChild}; planting the root at the local origin makes
\emph{any} node a valid sub-assembly root for the composite queries.}\label{lst:placement-resolvenodes}
\end{listing}

The resolve is not run per consumer. \cppinline|PartNode::ensurePlacementCache()|
(\src{core/model/PartsModel.cpp}{144}) runs \code{resolveNodes} and then \code{sweepOverlaps} once
per structural change --- node attached or detached, link edited, geometry mutated --- and caches
both the placement vector and the verdict on the node, \code{mutable} behind \code{const}
accessors. \cppinline|resolvedPlacements()| and \cppinline|placementDiagnostics()| serve the cached
values on every subsequent read; a mass change (a burning motor) never re-triggers them, because
geometry is invariant under a burn (\cref{sec:caching}).

\begin{designrule}[One resolve, many consumers]
Every consumer of absolute placement reads the same cached resolve through
\cppinline|resolvedPlacements()|: the composite mass/CM/inertia pass
(\src{core/model/PartsModel.cpp}{171}), the aero datum shift
(\src{core/model/PartsModel.cpp}{224}), the envelope sweep itself, the CLI's \code{checkdesign}
verdict (\src{cli/Repl.cpp}{969}), and the visualizer's mesh walk
(\src{visualizer/RocketMesh.cpp}{433}). Two independent resolvers could drift --- the simulator
flying one rocket while the renderer draws another --- so no consumer is permitted to compute its
own placement; disagreement is unrepresentable rather than tested for.
\end{designrule}

\subsection{Seam and envelope diagnostics}\label{sec:placement-diagnostics}

Diagnostics come in two granularities: a per-seam radial check for one attachment, and a whole-tree
envelope sweep over a resolved placement vector.

\paragraph{The radial seam check.}
\code{radialSeamCheck} (\src{core/model/parts/Placement.cpp}{57}) vets one parent--child edge,
dispatched on the seat kind exactly as \cref{tab:placement-seats} states: \seat{Abut} and
\seat{OnSurface} demand equal outer radii at the mated stations, \seat{NestInBore} demands the
child's OD fit within the parent's bore. It returns a located \code{OverlapDiagnostic} on
incompatibility and \code{nullopt} otherwise; its \code{zWorld} carries the parent-local seam
station, since no pose is in scope at seam granularity. Being free and pose-free, it is directly
exercisable by authoring flows and by the unit suite one seam at a time
(\srcf{core/model/tests/ResolverSweepTests.cpp}).

\paragraph{The envelope sweep.}
\code{sweepOverlaps} (\src{core/model/parts/Placement.cpp}{108}) checks the whole resolved tree for
physical impossibility: one part's material occupying another's. It works entirely on the
\code{Placed} vector --- \code{parentId} gives it the tree relations it needs, so it never touches
the ownership tree. For each part it builds the world axial interval
$[z_{\mathrm{fore}} - \mathrm{axialLength}, z_{\mathrm{fore}}]$, sorts intervals by aft station
($O(N \log N)$), and then tests each part (the \emph{offender}) against the part whose interior it
occupies (the \emph{host}) at a finite set of stations.

\begin{keyinsight}[Exact, not sampled]
The test stations are \emph{feature breakpoints}: the offender's own endpoints plus every other
interval endpoint falling within its span. For the stack profiles --- a tube constant, a cone
linear, a fin set's body disc constant --- the radius is monotone between breakpoints, so the
extremum of offender radius against host capacity is attained at a breakpoint and the sweep is
exact: no sampling density to tune, no thin intrusion slipping between samples. The one profile
outside this argument is \code{HollowSphere}, whose silhouette peaks at its equator between its
own endpoints; an intrusion carried only by that bulge, with no breakpoint inside the sphere's
span, is not sampled there.
\end{keyinsight}

At each breakpoint the host is chosen as the \emph{smallest-id covering interval}, after three
exclusions: the offender itself, its descendants (a child legitimately nests inside its parent's
envelope), and its direct seat parent (that relationship belongs to the seam check). Taking the
smallest id makes the report deterministic when several intervals cover the same station, so a
violation is always attributed to the same host regardless of sort incidentals. The capacity the
offender must fit within comes from \code{Part::innerCapacityAt}
(\src{core/model/parts/Part.cpp}{53}), which encodes the \emph{solid-host rule}: a bored part is
bounded by its inner wall, a solid part by its outer skin. Branching once inside the part keeps the
sweep free of solidity special cases --- the same comparison flags both a rod jammed through a
motor mount and a fin can poking out of a solid boattail.

Two further host rules are guards in the sweep loop itself (\cref{lst:placement-hostguards}):

\begin{listing}[htbp]
\begin{minted}{cpp}
if(cand.zFore - cand.zAft <= tol) { continue; }      // zero-span node (a Motor): no interior, hosts nothing
// ...
const double rOff = off.part->radiusOuterAt(zWorld - off.originZ);
const double cap  = host->part->innerCapacityAt(zWorld - host->originZ);

// A hollow host's bore capacity gates only an offender nested within its outer envelope. Once
// the offender's outer radius reaches the host skin (rOff >= hostOuter) it sits on or outside
// that skin -- e.g. a fin set's body disc on a co-radial aft coupler -- so a bore "intrusion"
// there is a modelling artifact, not a collision. A solid host has cap == hostOuter, so this
// guard never trips and poke-through still flags.
const double hostOuter = host->part->radiusOuterAt(zWorld - host->originZ);
const bool   boreHost  = cap < hostOuter - tol;        // hollow: capacity is the bore, not the skin
if(boreHost && rOff >= hostOuter - tol) { continue; }  // offender lies on/outside the host skin
\end{minted}
\caption{The zero-span-host rule and the bore-host skin guard inside the sweep
(\srccap{core/model/parts/Placement.cpp}{183}, elided).}\label{lst:placement-hostguards}
\end{listing}

\begin{designrule}[A zero-span node is mass, not envelope]
A node with no axial extent has no interior and never hosts. The canonical case is \code{Motor}
(\cref{sec:motor}), which keeps the base \code{getLength()} of $0$: it contributes mass, CM shift,
and inertia to the composite, but it is a point in the placement geometry, and a point cannot be
intruded into. Without this rule a zero-length interval would cover exactly its own station and, as
a low-id ``host'' of zero capacity, flag every co-located part as an offender.
\end{designrule}

The bore-host skin guard resolves the complementary edge case: a hollow host's bore capacity is
only meaningful for an offender \emph{inside} its outer envelope. An offender whose outer radius
reaches the host's skin sits on or outside that skin --- a fin set's body disc wrapped around a
co-radial coupler is the canonical case --- so comparing it against the bore would manufacture an
intrusion where none physically exists. Because a solid host's capacity equals its skin, the guard
never trips for solids and genuine poke-through still flags. Finally, the sweep reports only the
worst (deepest) violation per offender/host pair, so one misplaced part yields one located
diagnostic rather than a breakpoint-by-breakpoint flood.

\subsection{The diagnostics gate}\label{sec:placement-gate}

The sweep's verdict is cached beside the resolve it judged --- \cppinline|verdict_| is written in
the same \cppinline|ensurePlacementCache()| pass that fills \cppinline|resolved_|
(\src{core/model/PartsModel.cpp}{144}) --- so the verdict and the geometry it describes can never
be observed out of step, and no consumer can disagree with another about whether the design is
sound. \Cref{fig:placement-gate} traces the full path from a structural edit to the gated
consumers.

\tikzfig{activity-resolve-sweep}{From structural edit to verdict: the edit dirties the placement
chain up to the root; the next read resolves once, sweeps once, and caches both; every consumer
then sees the same placements and the same verdict.}{fig:placement-gate}

\begin{hardfail}[A self-intersecting design does not fly]
When the cached verdict is not \ok, \cppinline|computeCompositeAt| throws
\cppinline|std::runtime_error| carrying the first diagnostic's located message
(\src{core/model/PartsModel.cpp}{184}) --- and because cache stamps are written only after a
successful build (\cref{sec:caching}), the throw repeats on \emph{every} composite read rather than
being swallowed by the cache. A design whose geometry is physically impossible produces a hard,
located failure, never a silently wrong mass, CM, or inertia tensor fed to the integrator.
\end{hardfail}

Consumers differ only in how they surface the shared verdict. The flight path refuses to solve, as
above (\cref{sec:flight}). The CLI's \code{checkdesign} command (\src{cli/Repl.cpp}{969}) prints
each diagnostic's message behind an \code{ERR} verdict, and layers one additional check the sweep
deliberately does not make: an axial-coverage scan for air gaps, since a positive standoff is legal
in a \code{StationLink} but leaves parts floating apart --- a question of authored intent, not of
physical impossibility (\cref{sec:cli}). The visualizer (\src{visualizer/RocketMesh.cpp}{433})
keeps rendering --- a picture of the broken design is exactly what a user fixing it needs --- but
draws each offender in an error color and logs the same messages, so a headless caller sees the
same signal the simulator throws on. One resolve, one verdict, three presentations.
